\documentclass[conference]{IEEEtran}

\usepackage[T1]{fontenc}
\usepackage{cite}
\usepackage{url}
\usepackage{array}

\title{A Modular Scanner for Security Headers, Version Disclosure, Directory Listings, SQL Injection, and Reflected XSS}

\author{
\IEEEauthorblockN{Ayush Saraf}
\IEEEauthorblockA{Vellore Institute of Technology\\
ayush.saraf2023@vitstudent.ac.in}
\and
\IEEEauthorblockN{Shlok Lobhas Ghadekar}
\IEEEauthorblockA{Vellore Institute of Technology\\
shlok.lobhas2023@vitstudent.ac.in}
\and
\IEEEauthorblockN{Ankush}
\IEEEauthorblockA{Vellore Institute of Technology\\
ankush.2023@vitstudent.ac.in}
}

\begin{document}
\maketitle

\begin{abstract}
This paper presents a small Flask web application that runs five HTTP-based
security checks against a URL entered by the user. The checks look for missing
or weak response headers, exact software versions in headers, browsable
directories, database errors caused by form input, and unescaped reflected
markup. SQL injection and XSS checks share an HTML form parser; all five use a
common HTTP client. In the project’s manual verification record, the local
practice target produced five findings at \texttt{/login}. A scan of its root
page produced three findings; the SQL injection and XSS checks had no form to
test there. The repository documents 45 automated tests for positive, negative,
and error cases. These results show that the scanner works on its included
practice target and fixtures. They do not establish accuracy on other sites.
The scanner does not crawl, authenticate, run JavaScript, detect blind SQL
injection, or test stored and DOM-based XSS.
\end{abstract}

\begin{IEEEkeywords}
web application security, vulnerability scanning, SQL injection, reflected
cross-site scripting, HTTP security headers, security misconfiguration
\end{IEEEkeywords}

\section{Introduction}
Web applications expose endpoints that accept user input and return content to
browsers. A small scanner can help a student see how a request becomes a test
and how a response becomes a finding. This project focuses on five such tests:
three inspect headers or common paths, while two submit payloads through forms
found in the returned HTML. The checks address a narrow part of the risks
described in OWASP guidance for injection and security misconfiguration
\cite{owaspTop10,wstgSQLi,wstgXSS}; they are not a substitute for the broader
testing procedures in that guidance.

The scanner is a local Flask application. It accepts one HTTP or HTTPS URL and
executes five checks against the supplied page and a small set of common paths.
Two checks submit data to forms found in the returned HTML. Three inspect
response headers or bodies. Results appear in a report that gives the evidence,
severity, risk, and a suggested fix. A separate practice application supplies
known vulnerable behaviors so the end-to-end workflow can be tested without
relying on a changing external target.

The project was shaped by two practical questions: can five simple rules find
known behaviors in a deliberately vulnerable local application, and can
regression tests guard against specific mistakes in those rules? The
implementation brings the checks together through a coordinator, a shared
HTTP client, and a shared parser for HTML forms. A separate local target makes
the demonstration repeatable. The evaluation uses that target and the test
fixtures only; it does not estimate detection accuracy on arbitrary websites.

The remainder of the paper describes related testing guidance, the scanner
architecture, each detection rule, the development defects that informed the
regression tests, and the results documented for the practice target. It ends
with limitations and practical extensions.

\section{Background and Related Guidance}
The OWASP Top 10 groups injection and security misconfiguration among major
web-application risk categories \cite{owaspTop10}. The OWASP Web Security
Testing Guide describes testing approaches for SQL injection and reflected XSS
\cite{wstgSQLi,wstgXSS}. This project implements simplified response-based
probes inspired by these categories; it does not implement the full testing
procedures in that guide.

Alazmi and Conte de Leon reviewed published work on web application
vulnerability scanners and found that quantitative effectiveness evaluations
were limited and used differing approaches \cite{alazmi}. Their review
underscores why a scanner demonstration should not be confused with a
comparative performance study. Ibrahim and Rosli compare three tools for SQL
injection detection on predefined applications and report accuracy and response
time measures \cite{ibrahim}. The present project follows a smaller,
development-oriented approach: it checks known behaviors on its own practice
target and regression fixtures, without comparing scanner products or claiming
general detection rates.

HTTP response headers can communicate browser-enforced policy. Content Security
Policy (CSP) defines restrictions on resources a page may load or execute
\cite{csp3}. HTTP Strict Transport Security (HSTS) lets a site direct user
agents to use secure connections for a host \cite{rfc6797}. Other headers,
including \texttt{X-Frame-Options} and
\texttt{X-Content-Type-Options}, provide additional browser-side protections
when configured appropriately \cite{owaspHeaders}. The scanner treats the
presence and selected values of these headers as configuration signals; it
does not prove that the application's complete security policy is effective.

\section{System Architecture}
The system has four layers, shown in Fig.~\ref{fig:architecture}. The browser submits a target URL to the Flask
application in \texttt{app.py}. The application validates the URL scheme and
hostname shape, then calls \texttt{run\_all\_checks} in
\texttt{scanner/\_\_init\_\_.py}. The coordinator executes the registered
check modules in sequence and catches exceptions independently, producing an
\texttt{error} finding if one check cannot complete. The results are summarized
and rendered by the report template. An export route sanitizes posted finding
fields and returns a standalone HTML report with embedded CSS.

All checks use the shared HTTP module, \texttt{scanner/http.py}, which sets a
request timeout, a scanner user-agent, and a response-body size cap. Form-based
checks use \texttt{scanner/forms.py} to parse server-returned HTML with
BeautifulSoup and submit one field at a time. The form cache and HTTP-session
cookies are cleared before each scan. The check list is maintained in one
coordinator, while check-specific payloads, signatures, and finding logic stay
inside their respective modules.

\begin{figure}[!t]
\centering
\footnotesize
\begin{tabular}{c}
Browser submits URL\\
$\downarrow$\\
Flask input validation and scan route\\
$\downarrow$\\
Scanner coordinator: five independent checks\\
$\swarrow$ \hspace{0.2em} Shared HTTP layer \hspace{0.2em} $\searrow$\\
Target responses \hspace{1em} Shared HTML form parser\\
$\downarrow$\\
Findings, summary, report page, and HTML export
\end{tabular}
\caption{High-level scan and reporting flow.}
\label{fig:architecture}
\end{figure}

\section{Implementation Refinements}
Several regression cases in the repository correspond to concrete defects in
earlier implementations. They are useful to explain because they show how a
response-based scanner can report an incorrect result even when its payload
is reasonable.

\subsection{Database-Specific Error Signatures}
The SQL injection detector originally lacked signatures for the error wording
returned by the practice target's SQLite database. In that state, a quote
payload could break the query while the scanner reported no match. The current
signature set includes SQLite messages such as \texttt{unrecognized token} and
syntax errors near a token, alongside signatures for other database families.
The regression target emits the SQLite error and verifies that the SQL
injection check reports it. A second target safely escapes the submitted input
and is expected to pass.

\subsection{Path Resolution and Directory-Listing Classification}
Concatenating the submitted page URL with a probe path caused a target such as
\texttt{/login} to produce invalid probes beneath that page, such as
\texttt{/login/uploads/}. The shared \texttt{http.join} routine now removes
query strings and fragments and resolves the common probe paths against the
site's directory base. The directory check also requires both an HTTP 200
response and a listing signature in the body. This second condition matters
for single-page applications or catch-all routes that return their normal
shell with status 200 for every path. Regression tests cover a real listing, a
page-level target, a 403 response, and a 200-everywhere application shell.

\subsection{Reflected-XSS False-Negative Guard}
An earlier XSS condition required an escaped copy of the payload to be absent
from the entire response. That condition could cancel a real finding when a
page reflected the payload raw in its message and also placed an escaped copy
in a sticky form field. The current check looks for the submitted payload as a
raw substring; a correctly escaped response does not contain that raw markup.
Tests cover escaped reflection, raw reflection, raw reflection accompanied by
an escaped duplicate elsewhere, and a page with no forms. This rule reduces
one identified false-negative case, but it remains a heuristic: it does not
parse the browser execution context or execute the resulting page.

\subsection{Scan Isolation and Failure Reporting}
The request session and parsed-form cache are process-level objects. The
coordinator clears both at the beginning of each scan so cookies and cached
forms do not carry over between targets. It catches exceptions for each check
separately and returns an \texttt{error} finding, allowing the remaining checks
to complete. A regression case for a down target ensures connection failure
is reported as an error rather than a clean pass.

\section{Detection Methods}
Each check returns a dictionary containing a name, status, severity, detail,
risk, and suggested fix. Status is one of \texttt{pass}, \texttt{fail}, or
\texttt{error}. Severity is assigned using the project's check-specific rules;
it is not a calibrated measure of exploitability.

\subsection{Security Response Headers}
The header check sends a GET request and inspects the response for CSP, HSTS,
\texttt{X-Frame-Options}, and \texttt{X-Content-Type-Options}. It flags missing
headers and validates selected values: \texttt{X-Frame-Options} must be
\texttt{DENY} or \texttt{SAMEORIGIN}, while
\texttt{X-Content-Type-Options} must be \texttt{nosniff}. HSTS is skipped when
the final response uses HTTP because browsers do not apply the HSTS policy from
an insecure HTTP response \cite{rfc6797}. This is a header-level check and does
not evaluate the quality of a full CSP policy.

\subsection{Software-Version Disclosure}
The version check inspects \texttt{Server}, \texttt{X-Powered-By},
\texttt{X-AspNet-Version}, and \texttt{X-Generator}. A regular expression
searches for a product followed by a dotted version number, such as
\texttt{Apache/2.2.8}; a bare product name does not trigger the exact-version
finding. This check identifies selected header disclosures and does not
fingerprint software through other response features.

\subsection{Directory-Listing Exposure}
The directory check requests six common paths: \texttt{/admin/},
\texttt{/backup/}, \texttt{/uploads/}, \texttt{/files/}, \texttt{/images/},
and \texttt{/.git/}. A response is reported only when it has status 200 and
its first 4096 characters contain a known listing signature, such as
\texttt{Index of /} or \texttt{Parent Directory}. Redirects are not followed
for these probes so a redirect to a login page is not treated as an exposed
listing. The check does not discover arbitrary directories or filenames.

\subsection{Error-Based SQL Injection}
The shared form parser finds named text-like inputs in static HTML forms and
records each form's method and action. The SQL injection check submits four
basic strings to each injectable field and searches response text for database
error signatures associated with MySQL, Oracle, SQL Server, PostgreSQL, and
SQLite. A matching error produces a high-severity finding. The detector is
limited to visible errors; it does not infer boolean, time-based, or other blind
injection behavior.

\subsection{Reflected Cross-Site Scripting}
The XSS check submits three markup-bearing payload forms to each text field,
using a marker to make the test reflection recognizable. It reports a finding
when the submitted payload appears in the response body with its markup intact.
This is a lightweight reflection heuristic, not a browser-based execution
test. It therefore does not establish that the payload executes in every
browser context and does not inspect stored content or client-side DOM flows.

\section{Controlled Demonstration and Evaluation}
\subsection{Test Environment and Procedure}
The project includes a scanner app and a separate Flask practice target. The
scanner's local development server uses port 5000; the practice target uses
port 8000 and binds to the loopback interface. The practice target stores two
demonstration users in an in-memory SQLite database. It intentionally returns
database errors, reflects input in selected forms, omits selected response
headers, advertises sample version strings, and returns an index-like response
at \texttt{/uploads/}.

The manual verification record, dated July 20, 2026, scans the login page and
the root page as separate targets. The login page contains text fields and
exercises the form-based checks. The root page has no form, so its two form
check passes mean only that there was no input to submit on that page. The guide
also records HTTP responses inspected with \texttt{curl} and a manual
reproduction of the SQL injection behavior. Automated tests use temporary
HTTP-server fixtures for controlled positive, negative, and error responses.
The README lists 45 tests covering these behaviors and report export
\cite{projectReadme,manualVerification}.

\subsection{Observed Results}
The project verification record reports five failed checks and no passes for
\texttt{/login}. The findings are consistent with the practice target's
intentional configuration and route behavior. For \texttt{/}, three checks
fail and two form checks pass because no forms exist on that page. Those two
passes are scoped observations, not evidence that the entire practice site is
free of injection vulnerabilities.

\begin{table}[!t]
\caption{Documented Scan Outcomes}
\label{tab:results}
\centering
\footnotesize
\begin{tabular}{|p{0.18\columnwidth}|p{0.20\columnwidth}|p{0.50\columnwidth}|}
\hline
\textbf{Target} & \textbf{Summary} & \textbf{Interpretation} \\
\hline
\texttt{/login} & 5 fail, 0 pass & All five configured checks return a finding on the demonstration page or its common-path probes. \\
\hline
\texttt{/} & 3 fail, 2 pass & SQL injection and XSS pass because this page has no forms to submit to. \\
\hline
\end{tabular}
\end{table}

Table~\ref{tab:checks} connects the five reported findings with the evidence
recorded in the manual verification guide.
\begin{table}[!t]
\caption{Finding Evidence on the Practice Target}
\label{tab:checks}
\centering
\footnotesize
\begin{tabular}{|p{0.31\columnwidth}|p{0.57\columnwidth}|}
\hline
\textbf{Check} & \textbf{Documented evidence} \\
\hline
Security headers & CSP, X-Frame-Options, and X-Content-Type-Options are absent. \\
\hline
Version disclosure & Response headers identify Apache 2.2.8 and PHP 5.4.45. \\
\hline
Directory listing & \texttt{/uploads/} returns an index-like listing. \\
\hline
SQL injection & A quote in a login field causes a SQLite error; the manual guide records an authentication-bypass demonstration. \\
\hline
Reflected XSS & Test markup is returned raw in selected form responses. \\
\hline
\end{tabular}
\end{table}

\subsection{Regression Coverage and Interpretation}
The regression suite is designed to distinguish selected vulnerable behaviors
from common benign or failure cases. Examples include an escaped reflection
that should not trigger XSS, a page that returns status 200 for every path but
does not contain a directory listing, an exact server version versus a bare
product name, a weak header value, and an unreachable target. The suite also
checks that malformed export input produces a report rather than terminating
the request with an unhandled exception.

These tests are targeted regression checks, not a representative benchmark.
The practice target intentionally includes at least one example of each
configured indicator, and the tests cover a limited collection of safe and
edge-case responses. No labeled dataset was used to calculate precision,
recall, false-positive rate, or false-negative rate. No scan-time measurements
were collected. Accordingly, the result supports a claim of functional
operation on the included fixtures only; it cannot support a general claim of
effectiveness on arbitrary sites.

\section{Discussion}
The main engineering choice is to favor small, inspectable rules over broad
automatic discovery. This makes each report explainable: a reader can trace a
finding to a missing header, a matched version pattern, a directory-listing
signature, a database error, or a raw reflected payload. The same simplicity
limits the scanner. A rule can only identify cases represented by its selected
paths, payloads, response signatures, and form parser. The reported severity
is a static category attached to the rule, not a contextual risk calculation.

The regression history also shows why scanner output needs carefully scoped
wording. A response code alone was insufficient for directory-listing
detection, and searching for any echo of user input was insufficient for XSS.
Requiring a body signature and matching the raw submitted XSS string improved
these specific cases. They do not eliminate all false reports or missed
findings. A browser-based assessment, more varied targets, and independent
manual review would be needed to evaluate broader behavior.

The current implementation uses an eight-second HTTP timeout and caps each
response body at two megabytes. It disables TLS certificate verification to
support self-signed laboratory targets. This improves compatibility with
practice environments but removes certificate identity validation from HTTPS
requests. The application accepts any syntactically valid HTTP or HTTPS target;
the scanner binds to loopback, but the target itself is not technically
restricted. Operators must have permission to test any target they enter.

\section{Limitations, Safety, and Future Work}
The scanner accepts one HTTP or HTTPS URL and does not crawl linked pages,
authenticate, or execute JavaScript. Consequently, pages behind a login and
client-rendered forms are generally outside its coverage. A \texttt{pass}
means only that the implemented heuristic did not identify a finding for the
tested URL. It must not be interpreted as proof that the site is secure.

The HTTP client currently disables TLS certificate verification to support
self-signed certificates often used in local practice environments. The
scanner sends active payloads, and the target URL is not restricted to
loopback or private addresses. Testing should therefore be limited to systems
for which the operator has explicit authorization. The Flask application is
configured to bind to loopback for local use; it is not designed as a public
scanning service.

Possible extensions include a configurable crawl scope, optional
authentication, browser automation for JavaScript-rendered pages, broader
directory discovery, and additional SQL injection and XSS analysis. A separate
research evaluation could use a labeled set of vulnerable and hardened targets
to report precision, recall, and false-positive rates. These are extensions
beyond the current prototype's defined scope.

\section{Conclusion}
This paper described a modular, rule-based web vulnerability scanner with a
Flask interface, five basic checks, a shared HTTP and form layer, a downloadable
report, and a local vulnerable practice target. Its project documentation
reports successful operation against that fixture and includes automated
tests. The implementation is useful for demonstrating how response-based
checks can be organized and explained, while its narrow coverage and heuristic
matching limit conclusions about real-world security. Broader testing,
authentication and JavaScript support, TLS verification, and a reproducible
labeled evaluation are appropriate next steps.

\begin{thebibliography}{00}

\bibitem{owaspTop10}
OWASP Foundation, ``OWASP Top 10:2021,'' 2021. [Online]. Available:
\url{https://owasp.org/Top10/2021/}. [Accessed: Oct. 24, 2026].

\bibitem{wstgSQLi}
OWASP Foundation, ``SQL Injection,'' OWASP Web Security Testing Guide, latest
online edition. [Online]. Available:
\url{https://wstg.owasp.org/latest/4-Web_Application_Security_Testing/07-Injection/05-SQL_Injection/}.
 [Accessed: Oct. 24, 2026].

\bibitem{wstgXSS}
OWASP Foundation, ``Reflected Cross Site Scripting,'' OWASP Web Security
Testing Guide, latest online edition. [Online]. Available:
\url{https://wstg.owasp.org/latest/4-Web_Application_Security_Testing/07-Injection/01-Reflected_Cross_Site_Scripting/}.
 [Accessed: Oct. 24, 2026].

\bibitem{csp3}
M. West and A. Sartori, eds., ``Content Security Policy Level 3,'' W3C,
Working Draft, Sep. 16, 2026. [Online]. Available:
\url{https://www.w3.org/TR/CSP/}. [Accessed: Oct. 24, 2026].

\bibitem{rfc6797}
J. Hodges, C. Jackson, and A. Barth, ``HTTP Strict Transport Security
(HSTS),'' RFC 6797, Nov. 2012. [Online]. Available:
\url{https://www.rfc-editor.org/rfc/rfc6797}.

\bibitem{owaspHeaders}
OWASP Foundation, ``HTTP Headers Cheat Sheet,'' OWASP Cheat Sheet Series.
[Online]. Available:
\url{https://cheatsheetseries.owasp.org/cheatsheets/HTTP_Headers_Cheat_Sheet.html}.
 [Accessed: Oct. 24, 2026].

\bibitem{alazmi}
S. Alazmi and D. Conte de Leon, ``A systematic literature review on the
characteristics and effectiveness of web application vulnerability scanners,''
\emph{IEEE Access}, vol. 10, pp. 33200--33219, 2022, doi:
10.1109/ACCESS.2022.3161522.

\bibitem{ibrahim}
R. Y. Ibrahim and M. M. Rosli, ``Evaluation of web application vulnerability
scanners using SQL injection attacks,'' in \emph{Proc. 2023 IEEE 8th Int. Conf.
Recent Advances and Innovations in Engineering (ICRAIE)}, Kuala Lumpur,
Malaysia, Dec. 2023, pp. 1--6, doi: 10.1109/ICRAIE59459.2023.10468295.

\bibitem{projectReadme}
``Lightweight Web Vulnerability Scanner,'' project README and source code,
unpublished software project, 2026.

\bibitem{manualVerification}
``Manual Verification,'' project documentation, unpublished software project,
Jul. 20, 2026.

\end{thebibliography}
\end{document}
